% Drop-in replacements and additions for the ICML draft.
% All numbers are measured; provenance is in docs/THESIS_AND_RESULTS.md and the
% run directory named in each paragraph. Nothing here is a placeholder.

% =====================================================================
% 1. CONTRIBUTION 1, restated. The draft says the default fails "whenever
%    the model is accurate and the policy is variable". That is the symptom.
%    The cause is structural, and saying it turns a finding into a rule.
% =====================================================================
\paragraph{Contribution 1 (revised).}
\textbf{A structural account of when the default is wrong.}
Matched-pair residuals $c_i - g_i$ cancel any scatter common to both terms.
When the controller holds a \emph{per-instance} quantity $g_i$ at the limit,
as a control barrier function does pointwise, that cancellation is exactly
right. A Lagrangian dual in a CMDP holds an \emph{expectation}
$\mathbb{E}[c]$, so $g_i$ is not the controlled quantity and the per-instance
scatter that breaches $d$ cancels out of the calibration set. The default is
sound in the setting it was developed for and unsound here, and the
distinction can be stated in advance rather than discovered.

% =====================================================================
% 2. SECTION 4.4 replacement. One quantile, one assumption, no bias term,
%    and the matched-pair recipe falls out as the special case.
% =====================================================================
\subsection{Our margin: conformalise against the controlled quantity}
\label{sec:residual}

Let $\hat g$ denote the quantity the dual actually holds at the limit: the
model's estimate of $J_C$ for the current policy. The probe returns realised
episode costs $c_1,\dots,c_m$ under $F$ for that same policy. We calibrate on
\begin{equation}
  s_i \;=\; c_i - \hat g,
  \qquad
  q \;=\; s_{(\lceil (n+1)(1-\delta)\rceil)},
  \qquad
  d_{\mathrm{eff}} \;=\; d - q ,
  \label{eq:residual}
\end{equation}
pooling $s_i$ across probes. A single order statistic absorbs both mechanisms
of Section~4.1: systematic bias shifts every $s_i$, and episode spread
disperses them. No separate bias term is required.

\begin{assumption}[Exchangeable residuals]\label{as:exch}
A new evaluation episode has cost $c = \hat g_{\mathrm{eval}} + s$, where $s$
is exchangeable with the calibration residuals $s_1,\dots,s_n$.
\end{assumption}

\begin{proposition}\label{prop:cov}
Under Assumption~\ref{as:exch}, if the dual holds the model's estimate at the
tightened limit, $\hat g_{\mathrm{eval}} \le d_{\mathrm{eff}}$, then
$\mathbb{P}[c > d] \le \delta$, marginally over the calibration draw.
\end{proposition}

\begin{proof}
$\hat g_{\mathrm{eval}} \le d - q$ implies $\{c > d\} \subseteq \{s > q\}$. By
exchangeability the rank of $s$ among $\{s_1,\dots,s_n,s\}$ is uniform, so
$\mathbb{P}[s \le s_{(k)}] \ge k/(n+1) \ge 1-\delta$ for
$k = \lceil (n+1)(1-\delta)\rceil$ \citep{lei2018splitconformal}.
\end{proof}

\begin{remark}[The default is the special case]
If the controller holds a per-instance quantity, so that $\hat g$ is replaced
by $g_i$ for each calibration episode, then $s_i = c_i - g_i$ is precisely the
matched-pair residual of Section~4.3 and Proposition~\ref{prop:cov} recovers
the standard guarantee. The two recipes are not rivals; they answer the same
question for different controllers, and the CMDP dual is the case where they
part company.
\end{remark}

\begin{proposition}[When a margin is usable at all]\label{prop:feasible}
Let $s$ be the span of the constrained quantity between the constraint-optimal
and objective-optimal policies, $b$ the model-reality bias,
$\sigma_{\mathrm{agg}}$ the spread of the quantity violation is declared on,
and place the limit at $d = d_{\mathrm{best}} + f s$. Since the margin is a
quantile of the residuals rather than the bias alone, $d_{\mathrm{eff}}$ is
reachable by the policy only if
\begin{equation}
  b + z_\delta\,\sigma_{\mathrm{agg}} \;<\; f\,s .
  \label{eq:feasible}
\end{equation}
Otherwise the constraint is unachievable however well the margin is calibrated,
and a violation rate cannot distinguish that case from a miscalibrated one.
\end{proposition}

\begin{figure}[t]
\centering
\includegraphics[width=\linewidth]{figures/paper/fig3_precondition.pdf}
\caption{Equation~\eqref{eq:feasible} on the wildfire task. Points below the
diagonal admit a margin. Four settings of crew budget and model capacity
(crosses) never reach it: the bias runs $2.0$ to $4.9\times$ the span the
planner controls. With the discrepancy dialled instead (circles) the boundary
is straddled, which is what a demonstration of the condition requires.}
\label{fig:precondition}
\end{figure}

\begin{remark}[$f$ is the trade-off, not a free parameter]
It sets how tightly the constraint binds and how much room a margin has, in
opposite directions. Reporting $f$, $b$ and $s$ together is the minimum for a
reader to judge whether a margin could have worked at all.
\end{remark}

\begin{remark}[Weaker than the draft's assumption]
Equation~\eqref{eq:residual} requires one exchangeable sample.  The
deviations-plus-bias construction requires an exchangeable sample
\emph{and} a bias estimate that carries no coverage statement of its own.
Assumption~\ref{as:exch} is still strong across a learning run, since the
policy moves between probes; this is the drift adaptive conformal inference
was designed for \citep{gibbs2021aci}, and we report the empirical rate rather
than claim the assumption holds.
\end{remark}

% =====================================================================
% 2b. The calibration curve. A conformal method is finally judged on whether
%     the rate it states is the rate it delivers, and the other experiments
%     all fix delta = 0.1.
% =====================================================================
\subsection{Does the stated level mean anything?}
\label{sec:calibration}

A conformal method is judged on whether the level it states is the level it
delivers. The experiments above fix $\delta = 0.1$; here $\delta$ is swept, in
two regimes, with everything else controlled (Figure~\ref{fig:calibration}).

\begin{figure}[t]
\centering
\includegraphics[width=\linewidth]{figures/paper/fig0_calibration.pdf}
\caption{Realised violation rate against the stated level, $600$ calibration
draws per point. Left, where the model's per-instance error dominates
($\rho=0.25$), the default is merely over-conservative. Right, where the
policy's spread dominates ($\rho=4$), its realised rate is nearly independent
of the level requested: $0.30$ when asked for $0.02$, $0.47$ when asked for
$0.40$. Ours tracks the diagonal in both.}
\label{fig:calibration}
\end{figure}

The right panel is the substantive one. The default does not merely undercover;
its realised rate barely responds to $\delta$ at all, so the number it states
carries almost no information about its behaviour. A practitioner who tightens
$\delta$ from $0.1$ to $0.02$ would obtain $0.30$ instead of $0.37$, having
asked for a fivefold reduction.

% =====================================================================
% 3. NEW SUBSECTION for Section 5. This is the head-to-head the draft
%    deferred to future work, and without it a reviewer can dismiss the
%    18.2% result as a mis-tuned implementation.
% =====================================================================
\subsection{Which recipe, and when}
\label{sec:crossover}

Section~5 showed the default failing where the surrogate is accurate. If that
is structural rather than an artefact of our implementation, degrading the
surrogate should hand the advantage back. We sweep surrogate fidelity by
training budget---$1$ epoch on $20$ trajectories through $20$ epochs on
$256$---holding everything else fixed, and run all four margins at each level
over five seeds. Held-out relative $L_2$ ranges from $0.82$ to $0.011$.

Write $\rho = \sigma_{\mathrm{ep}} / \varepsilon$ for the policy's episode-cost
spread over the surrogate's per-instance cost error. Coverage at level $\delta$
requires a margin of $b + z_\delta \sigma_{\mathrm{ep}}$; the matched-pair
quantile supplies approximately $b + z_\delta \varepsilon$; the bias cancels.
The default therefore covers exactly when $\varepsilon \ge \sigma_{\mathrm{ep}}$,
with the crossing at $\rho = 1$.

\begin{table}[t]
\caption{Violation rate by regime, $25$ runs, $41$ evaluations each,
$\delta = 0.1$. Where the model's per-instance error dominates the default is
correct; where the policy's spread dominates it coincides with applying no
margin at all.}
\label{tab:crossover}
\centering\small
\begin{tabular}{lcccc}
\toprule
& none & model err. & deviations & \textbf{ours} \\
\midrule
$\rho<8$ \small(17 runs) & $0.116$ & $\mathbf{0.072}$ & $0.060$ & $0.060$ \\
$\rho\ge8$ \small(8 runs) & $0.140$ & $\mathbf{0.140}$ & $0.067$ & $0.073$ \\
\bottomrule
\end{tabular}
\end{table}

\begin{figure}[t]
\centering
\includegraphics[width=\linewidth]{figures/paper/fig1_mechanism.pdf}
\caption{(a) Coverage in the controlled setting: the default fails once
$\rho>1$. (b) On the reaction--diffusion front, the margin each recipe
produces against the $\rho$ measured in that run. The default's margin
collapses by a factor of $48$ as the surrogate improves, from $0.385$ to
$0.008$ against a limit that stays at $3.26$; ours is flat
($0.616 \pm 0.143$) because it tracks the policy's spread, which does not
improve with the model.}
\label{fig:mechanism}
\end{figure}

Figure~\ref{fig:mechanism}(b) shows the cause directly: the default's margin
collapses $48\times$ across the fidelity ladder while ours is flat.
Table~\ref{tab:crossover} reports the consequence. Below the crossing the default
holds $\delta$ and is statistically indistinguishable from ours ($z = 0.86$);
the recipe is not wrong in general, which is what makes this a condition of use
rather than a criticism. Above it the default returns $0.140$ against $0.140$
for no margin at all---the same value to three decimals---and exceeds the level
it states, while ours holds $0.073$ ($z = 2.78$). Ours improves on no margin in
both regimes ($z = 3.68$ and $z = 2.78$).

\paragraph{What the margin costs.}
Figure~\ref{fig:tradeoff} answers the first objection to any conservative
method. Paired against the same run's unconstrained arm, ours removes $4.8$ points of
violation for $1.8\%$ of return; the deviations-plus-bias form, $4.2$ points
for $1.5\%$. Safety here is not bought by refusing to act.

\begin{figure}[t]
\centering
\includegraphics[width=\linewidth]{figures/paper/fig5_tradeoff.pdf}
\caption{Change in violation rate and return against the same run's
unconstrained arm, means with bootstrap intervals over $25$ runs. Individual
runs are not shown: at $41$ evaluations each the per-run violation change is
quantised, so a scatter of the raw pairs shows binning rather than a
relationship.}
\label{fig:tradeoff}
\end{figure}

The coincidence at high $\rho$ is the sharpest form of the result. The margin is
not merely loose: the quantity it conformalises has stopped carrying
information about what breaches the limit, so tightening by it changes nothing.

\paragraph{Against a model-based baseline.}
Our other safe-RL baselines are model-free, which makes a sample-efficiency
comparison uninformative for a model-based planner. CAP \citep{ma2022cap} is
the closest method that also plans in a learned model and corrects for its
error, penalising by ensemble disagreement $c = \bar{c} + k\,\mathrm{sd}(c)$
over five surrogates with $k$ adapted from real data. We transplant that
penalty onto our planner at an identical probe budget, so only the correction
mechanism differs (Appendix~\ref{app:cap}).

\begin{table}[t]
\caption{Matched real-sample budget ($6{,}400$ training transitions plus
$1{,}920$ probe), \textsc{rdf}, full-fidelity surrogate. CAP is
indistinguishable from applying no margin at all.}
\label{tab:cap}
\centering\small
\begin{tabular}{lcc}
\toprule
method & violation & seeds \\
\midrule
no margin & $0.132$ & 5 \\
model-error conformal & $0.142$ & 5 \\
CAP (ensemble penalty) & $0.120$ & 12 \\
\textbf{ours} & $\mathbf{0.078}$ & 10 \\
\midrule
\multicolumn{3}{l}{\small CAP vs.\ no margin\quad $z=-0.43$} \\
\multicolumn{3}{l}{\small ours vs.\ no margin\quad $z=+3.68$} \\
\multicolumn{3}{l}{\small \textbf{ours vs.\ CAP\quad $z=+2.08$}} \\
\bottomrule
\end{tabular}
\end{table}

Table~\ref{tab:cap} reports the comparison. CAP's adapted $k$ settled between
$0.50$ and $1.46$, so its mechanism was active rather than inert; ensemble
disagreement is simply not the quantity that breaches this limit, for the same
structural reason model error is not. CAP matches on real samples but uses five
times the training compute, the ensemble being five models fitted to the same
data; a compute-matched comparison would be less generous to it, and a
sample-matched one is the right choice for a claim about constraint
satisfaction rather than efficiency.

\paragraph{The general form.}
Two mechanistically different uncertainty corrections---a conformal quantile
over model error, and an ensemble-variance penalty---both reduce to doing
nothing on an expectation constraint, while conformalising against the quantity
the dual controls works. Both estimate \emph{how wrong the model is}, and under
an expectation constraint that is not what carries the policy over the limit.

% =====================================================================
% 4. NEW SUBSECTION for Section 6. The draft states outright that Section 6
%    validates the planner and not the margin. This closes that gap.
% =====================================================================
\subsection{The margin on real fire dynamics}
\label{sec:ndwsmargin}

Section~6 enforced the crew budget by projection, so it exercised the planner
and not the margin. We close that gap without weakening the limitation of
Section~7.

The model/reality gap cannot be built from a second trained model on this task.
Over four settings of crew budget and model capacity the bias ran $2.0$ to
$4.9\times$ the entire span the planner can move the constrained quantity over,
so the constraint was unachievable and no margin could be assessed. More crew
made it worse, both the objective-optimal and constraint-optimal plans
converging on the same floor. This is
sim2sim: $F$ is a model, and no observational record contains a fire where a
firebreak was cut on our instruction. What it supplies is a model/reality cost
gap that nobody tuned, over dynamics fitted to real fires rather than to a PDE
we wrote.

The crew budget remains a hard projection. The \emph{constraint} is an
expectation, as in any CMDP, and conflicts with the objective: we minimise
total burned area subject to population-weighted burn staying below $d$, which
bind whenever sparing a large empty area costs a small populated one. The
limit is calibrated between what the objective-optimal planner ($\lambda=0$)
leaves on the populated area and what a town-focused planner achieves, so
meeting it costs a real share of the objective. A single multiplier, shared
across fires and driven by the batch mean under $G$, prices the constraint
while each fire is still planned at decision time.

Violation is measured per evaluation batch, because the constrained quantity
is a mean. This is exactly the regime Section~\ref{sec:residual} predicts the
matched-pair recipe cannot cover: the across-fire scatter that pushes batch
means over the limit appears in both $c_i$ and $g_i$ and cancels.

\paragraph{Outcome: the task fails the precondition.}
Run either side of the threshold, the margin exceeded the feasible headroom at
every perturbation: at the smallest, headroom $0.0034$ against a required
margin of $0.0164$; at the largest, $0.0062$ against $0.0877$, which is $93\%$
of the limit itself. Violation stayed at $0.65$ to $0.80$ for every arm
including ours, because $d_{\mathrm{eff}}$ sits below anything the planner can
reach.

This is a property of the task. The constrained quantity's batch-to-batch
spread is comparable to the planner's entire authority over it, so no margin
covering that spread fits the headroom---at any perturbation, with any limit
placement. The only remaining lever, averaging over larger evaluation batches,
would need $1{,}903$ fires per batch against the $1{,}500$ held out, and would
leave under one batch from which to estimate a $10\%$ rate. Miscalibration, an
under-converged dual and insufficient crew were each tested and none was the
cause.

We report this as a worked example of a system that fails
Proposition~\ref{prop:feasible} rather than omitting it. A reader deciding
whether the method applies to their problem is better served by a case that
fails the condition than by another that passes it.

% =====================================================================
% 5. SECTION 2 replacement paragraph. The draft mischaracterises Prinster
%    et al., which is the closest prior work to ours.
% =====================================================================
\paragraph{Conformal prediction for control (revised).}
Split conformal prediction \citep{vovk2005alrw,lei2018splitconformal,
angelopoulos2023gentle} gives finite-sample, distribution-free coverage under
exchangeability; adaptive conformal inference \citep{gibbs2021aci} relaxes
exchangeability under drift. In control, the dominant instantiation calibrates
a quantile of the error between the learned dynamics and the true system and
tightens a barrier or Lyapunov constraint by it
\citep{lindemann2023safeplanning,zhou2025acpsafe,
mirzaeedodangeh2026robustcbf}, with cumulative bounds for nonstationary
constraints \citep{tomashevskiy2026proactive}. A separate strand calibrates
against a reference \emph{policy} rather than a model
\citep{prinster2026cpc}, which shares our instinct that the calibration target
should be the controlled object; we differ in conformalising realised cost
against the quantity a Lagrangian dual holds, which yields a CMDP budget
rather than a behavioural-change budget. Our setting is a CMDP with an
episode-cost budget rather than CBF-constrained continuous control, and we
have not reproduced the failure in theirs.
